\documentclass[12pt,a4paper]{article}

\makeatletter
\usepackage[utf8]{inputenc}
\usepackage[T2B]{fontenc}
\usepackage[english,russian]{babel}

\usepackage{geometry}
\geometry{
  a4paper,
  left=2.54cm,
  right=2.54cm,
  top=2.54cm,
  bottom=2.54cm,
  headheight=0mm,
  headsep=0mm,
  footnotesep=0mm,
  footskip=0mm,
}

\usepackage{nopageno}

\usepackage{tempora}
\renewcommand{\ttdefault}{cmtt}

\usepackage{titlesec}
\titleformat*{\section}{\large\bfseries}
\titlespacing\section{0pt}{12pt plus 4pt minus 2pt}{0pt plus 2pt minus 2pt}
\setcounter{secnumdepth}{0}

\usepackage{enumitem}

\frenchspacing

\usepackage{amsmath}
\usepackage{amssymb}
\usepackage{hyperref}
\usepackage{longtable}
\usepackage[table]{xcolor}
\usepackage{array}
\usepackage{color}
\usepackage{xcolor}

\usepackage{hyperref}

\newcommand{\MYhref}[3][blue]{\href{#2}{\color{#1}{#3}}}%

\usepackage{listings}
\usepackage{alltt}
\usepackage{csquotes}
\DeclareQuoteStyle{russian}
	{\guillemotleft}{\guillemotright}[0.025em]
	{\quotedblbase}{\textquotedblleft}
\ExecuteQuoteOptions{style=russian}

\usepackage{graphicx}

\lstset{tabsize=2,
	breaklines,
	columns=fullflexible,
	flexiblecolumns,
	numbers=left,
	numberstyle={\footnotesize},
	extendedchars,
	inputencoding=utf8}

\def\@xobeysp{ }
\def\verbatim@processline{\hspace{1.2cm}\raggedright\the\verbatim@line\par}

\parindent=0pt
\parskip=3pt

\definecolor{lightgray}{gray}{0.9}

\renewcommand{\thesubsection}{\arabic{subsection}}

\lstdefinestyle{customc}{
  belowcaptionskip=1\baselineskip,
  breaklines=true,
  frame=L,
  xleftmargin=\parindent,
  language=C,
  showstringspaces=false,
  basicstyle=\footnotesize\ttfamily,
  keywordstyle=\bfseries\color{green!40!black},
  commentstyle=\itshape\color{gray},
  identifierstyle=\color{black},
  stringstyle=\color{blue},
}

\lstdefinestyle{customasm}{
  belowcaptionskip=1\baselineskip,
  frame=L,
  xleftmargin=\parindent,
  language=[x86masm]Assembler,
  basicstyle=\footnotesize\ttfamily,
  commentstyle=\itshape\color{purple!40!black},
}

\lstset{escapechar=@,style=customc}
\makeatother

% Style for terminal output and strace logs: no C highlighting, no numbers.
\lstdefinestyle{console}{
  language={},
  numbers=none,
  frame=single,
  breaklines=true,
  columns=fullflexible,
  keepspaces=true,
  basicstyle=\footnotesize\ttfamily,
  identifierstyle=\color{black},
}

% listings cannot read multibyte UTF-8: listingsutf8 recodes the source files
% (Cyrillic comments) to KOI8-R on the fly.
\usepackage{listingsutf8}
\lstset{inputencoding=utf8/koi8-r}

\hypersetup{hidelinks}

% Monospaced identifiers cannot be hyphenated; let TeX stretch the spaces
% instead of pushing them into the margin.
\setlength{\emergencystretch}{3em}

% Every line number of a source listing becomes a link target "<file>.<line>".
\newcommand{\lstlinetarget}[1]{\raisebox{\baselineskip}[0pt][0pt]{\hypertarget{\lstname.#1}{}}#1}
\lstset{numberstyle=\footnotesize\lstlinetarget}

% Code references: \deffn{key}{file}{line} declares where a definition starts,
% \fn{key} prints the key as a link, \coderef{key}{text} prints custom text.
% The line numbers below must match the sources in ../src.
\newcommand{\deffn}[3]{\expandafter\def\csname cref@#1\endcsname{#2.#3}}
\newcommand{\coderef}[2]{%
  \ifcsname cref@#1\endcsname
    \hyperlink{\csname cref@#1\endcsname}{\textcolor{blue}{\texttt{#2}}}%
  \else
    \errmessage{Unknown code reference: #1}%
  \fi}
\newcommand{\fn}[1]{\coderef{#1}{\detokenize{#1}}}

\deffn{master.c}{master.c}{1}
\deffn{slave.c}{slave.c}{1}
\deffn{common.h}{common.h}{1}
\deffn{common.c}{common.c}{1}
\deffn{Makefile}{Makefile}{1}

\deffn{kChildrenCount}{master.c}{19}
\deffn{kErrInput}{master.c}{22}
\deffn{kErrSetup}{master.c}{23}
\deffn{kErrChild1Write}{master.c}{24}
\deffn{kErrChild2Write}{master.c}{25}
\deffn{kErrChild1Status}{master.c}{26}
\deffn{kErrChild2Status}{master.c}{27}
\deffn{child_t}{master.c}{29}
\deffn{check_std_fds}{master.c}{39}
\deffn{close_fd}{master.c}{53}
\deffn{set_cloexec}{master.c}{60}
\deffn{open_pipe}{master.c}{68}
\deffn{build_slave_path}{master.c}{80}
\deffn{read_filename}{master.c}{99}
\deffn{open_output_file}{master.c}{134}
\deffn{is_same_regular_file}{master.c}{146}
\deffn{spawn_slave}{master.c}{156}
\deffn{wait_child}{master.c}{174}
\deffn{start_children}{master.c}{213}
\deffn{forward_lines}{master.c}{267}
\deffn{finish_children}{master.c}{316}
\deffn{master-main}{master.c}{332}

\deffn{is_vowel}{slave.c}{8}
\deffn{slave-main}{slave.c}{12}

\deffn{kBlockSize}{common.h}{8}
\deffn{kChildSetupFailed}{common.h}{13}
\deffn{kChildExecFailed}{common.h}{14}
\deffn{kSlaveReadFailed}{common.h}{19}
\deffn{kSlaveWriteFailed}{common.h}{20}
\deffn{read_retry}{common.c}{7}
\deffn{write_all}{common.c}{15}

\begin{document}

\begin{titlepage}
\begin{center}
    \includegraphics[width=\textwidth]{mai_header.jpg}

    \vspace{1cm}
    Министерство науки и высшего образования Российской Федерации\\
    Федеральное государственное бюджетное образовательное учреждение высшего образования\\
    \textbf{Московский авиационный институт}\\
    (Национальный исследовательский университет)

    \vspace{1.5cm}
    Институт № 8 \enquote{Компьютерные науки и прикладная математика}\\
    Кафедра 806

    \vspace{2cm}
    {\large\textbf{Отчёт по лабораторной работе №2}}\\[0.3cm]
    {\textbf{Дисциплина}}: \enquote{Операционные системы и архитектура компьютера}\\[0.3cm]
    {\textbf{Тема}}: \enquote{Работа с процессами ОС}\\[0.3cm]
    {\textbf{Вариант}}: 18
\end{center}

\vspace{2cm}
\begin{flushright}
    {\textbf{Выполнил:}}\\
    студент группы М8О-206БВ-25\\
    Ширанков М.С.\\

    {\textbf{Проверил:}}\\
    Ассистент Центра ТОП-ИТ\\
    Ефимов Д.А.
\end{flushright}

\vfill
\begin{center}
    Москва 2026
\end{center}
\end{titlepage}

\tableofcontents
\newpage

\section{Цель работы}
Приобрести практические навыки управления процессами в ОС семейства UNIX:
порождения дочерних процессов (\texttt{fork}), замены их образа другой программой
(\texttt{exec}), перенаправления стандартных потоков (\texttt{dup2}) и обмена
данными между процессами через неименованные каналы (\texttt{pipe}), а также
корректной обработки ошибок всех используемых системных вызовов.

\section{Задание}
Родительский процесс создаёт два дочерних процесса. Первой строкой пользователь
вводит в консоль родительского процесса имя файла, который будет открыт на запись
для \texttt{child1}, второй строкой "--- аналогичное имя для \texttt{child2}.
Родительский и дочерний процессы должны быть представлены разными программами.

Родительский процесс принимает от пользователя строки произвольной длины и
пересылает их в \texttt{pipe1} или \texttt{pipe2} в зависимости от правила
фильтрации. Процессы \texttt{child1} и \texttt{child2} обрабатывают строки и пишут
результаты своей работы в стандартный вывод. Все возможные ошибки должны
корректно обрабатываться.

\textbf{Правило фильтрации (вариант 18):} нечётные строки отправляются в
\texttt{pipe1}, чётные "--- в \texttt{pipe2}. Дочерние процессы удаляют из строк
все гласные.

\section{Ход работы}
Решение состоит из двух программ и общего модуля (полный исходный код
приведён в приложении~А):
\begin{itemize}
    \item \fn{master.c} "--- родительский процесс: читает имена файлов, создаёт
          каналы, порождает дочерние процессы и распределяет строки между ними;
    \item \fn{slave.c} "--- дочерняя программа: фильтр
          \texttt{stdin}~$\rightarrow$~\texttt{stdout}, удаляющий гласные
          \texttt{aeiouAEIOU};
    \item \fn{common.c}, \fn{common.h} "--- общие функции ввода-вывода
          (\fn{read_retry}, \fn{write_all}) и коды завершения.
\end{itemize}

Обе программы запускаются от одного и того же исполняемого файла \texttt{slave}:
\texttt{child1} и \texttt{child2} "--- это два экземпляра одной программы,
различающиеся только подключёнными к ним каналом и файлом. Сборка выполняется
по \fn{Makefile} из корня лабораторной работы:

\begin{lstlisting}[style=console]
make
./master
\end{lstlisting}

Ключевое решение: файл для дочернего процесса открывает родитель
(\fn{open_output_file}), а затем через \texttt{dup2} подставляет его дескриптор
на место стандартного вывода ребёнка (\fn{spawn_slave}). Поэтому дочерняя программа, как и требует задание, пишет результат
в \texttt{stdout} и ничего не знает о файлах. Ошибка открытия файла при этом
обнаруживается до порождения процессов.

\section{Используемые системные вызовы}
В программе используются следующие системные вызовы и функции стандартной
библиотеки, являющиеся тонкими обёртками над ними:

\begin{longtable}{|p{0.14\textwidth}|>{\raggedright\arraybackslash}p{0.47\textwidth}|>{\raggedright\arraybackslash}p{0.29\textwidth}|}
\hline
\rowcolor{lightgray}
\textbf{Вызов} & \textbf{Назначение в программе} & \textbf{Где в коде} \\
\hline
\endhead
\texttt{read} & чтение имён файлов (по одному байту) и строк (блоками по 4~КБ)
               из \texttt{stdin}; чтение дочерним процессом своего канала &
               \fn{read_retry}, \fn{read_filename}, \fn{forward_lines},
               \coderef{slave-main}{main} (slave) \\
\hline
\texttt{write} & пересылка строк в каналы; запись результата дочерним процессом
                в \texttt{stdout} & \fn{write_all} \\
\hline
\texttt{open} & открытие (создание, усечение) выходных файлов с флагом
               \texttt{O\_CLOEXEC} & \fn{open_output_file} \\
\hline
\texttt{fstat} & проверка того, что оба имени не указывают на один и тот же файл &
                \fn{is_same_regular_file} \\
\hline
\texttt{pipe} & создание двух неименованных каналов \texttt{pipe1}, \texttt{pipe2} &
               \fn{open_pipe} \\
\hline
\texttt{fcntl} & проверка открытости дескрипторов 0--2 (\texttt{F\_GETFD});
                установка флага \texttt{FD\_CLOEXEC} на концы каналов &
                \fn{check_std_fds}, \fn{set_cloexec} \\
\hline
\texttt{fork} & порождение дочернего процесса & \fn{spawn_slave} \\
\hline
\texttt{dup2} & перенаправление \texttt{stdin} ребёнка на канал, а \texttt{stdout}
               "--- на файл & \fn{spawn_slave} \\
\hline
\texttt{execv} & замена образа дочернего процесса программой \texttt{slave} &
                \fn{spawn_slave} \\
\hline
\texttt{sigaction} & игнорирование \texttt{SIGPIPE} в родителе: запись в канал без
                    читателя возвращает \texttt{EPIPE}, а не убивает процесс &
                    \fn{start_children} \\
\hline
\texttt{close} & закрытие ненужных концов каналов; закрытие пишущих концов для
                передачи детям признака конца данных (EOF) &
                \fn{close_fd}, \fn{finish_children} \\
\hline
\texttt{waitpid} & ожидание завершения детей и разбор их кода возврата &
                  \fn{wait_child} \\
\hline
\texttt{\_exit} & завершение ребёнка при ошибке \texttt{dup2}/\texttt{execv} без
                 сброса унаследованных буферов \texttt{stdio} & \fn{spawn_slave} \\*
\hline
\end{longtable}

\subsection{Подготовка}
Подготовку выполняет функция \fn{start_children}:
\begin{enumerate}
    \item \textbf{Проверка стандартных дескрипторов.} Если программа запущена с
          закрытым \texttt{stdin} (\texttt{./master <\&-}), то \texttt{pipe()}
          вернул бы свободный номер 0, и последующий \texttt{dup2} перепутал бы
          дескрипторы. Поэтому сначала \fn{check_std_fds} с помощью
          \texttt{fcntl(fd, F\_GETFD)} проверяет, что 0, 1 и 2 открыты.
    \item \textbf{Чтение имён файлов} (\fn{read_filename}). Имена читаются по одному байту до
          \texttt{\textbackslash n}. Так родитель забирает из \texttt{stdin} ровно
          две строки, а остальные данные остаются в ядре для этапа пересылки.
          Длина имени ограничена \texttt{PATH\_MAX}; завершающий
          \texttt{\textbackslash r} отбрасывается.
    \item \textbf{Открытие файлов} (\fn{open_output_file}) с флагами \texttt{O\_WRONLY},
          \texttt{O\_CREAT}, \texttt{O\_TRUNC} и \texttt{O\_CLOEXEC}. Затем
          \fn{is_same_regular_file} сравнивает пары
          \texttt{(st\_dev, st\_ino)}: два ребёнка с независимыми
          смещениями в одном обычном файле затирали бы данные друг друга.
    \item \textbf{Создание каналов} (\fn{open_pipe}) и установка \texttt{FD\_CLOEXEC}
          на оба их конца (\fn{set_cloexec}).
\end{enumerate}

\subsection{Порождение дочерних процессов}
Каждый ребёнок запускается функцией \fn{spawn_slave}. После \texttt{fork} ребёнок получает копию всех дескрипторов родителя.
\texttt{dup2} ставит нужный конец канала и файл на места 0 и 1, причём новый
дескриптор флага \texttt{FD\_CLOEXEC} не наследует. Все остальные копии ядро
закрывает автоматически при \texttt{execv}:

\begin{lstlisting}[style=console]
child1 after fork:   0=stdin 1=stdout 2=stderr 3=file1 4=file2 5=pipe1[0] 6=pipe1[1] 7=pipe2[0] 8=pipe2[1]
after dup2:          0=pipe1[0] 1=file1 2=stderr   3..8 are copies with CLOEXEC
after execv:         0=pipe1[0] 1=file1 2=stderr   (3..8 closed by the kernel)
\end{lstlisting}

Если бы \texttt{child1} сохранил копию \texttt{pipe2[1]},
\texttt{child2} не получил бы конец данных, пока жив \texttt{child1}. Путь к
программе \texttt{slave} строится от \texttt{argv[0]} (\fn{build_slave_path}),
поэтому \texttt{master} можно запускать из любого каталога. После каждого
\texttt{fork} родитель закрывает свои копии читающего конца канала и файла
(\fn{close_fd} в \fn{start_children}).

Обработчик \texttt{SIGPIPE} устанавливается после \texttt{fork}, так как
диспозиция \texttt{SIG\_IGN} наследуется через \texttt{exec}, а детям нужно
поведение по умолчанию.

\subsection{Пересылка строк}
Пересылку выполняет функция \fn{forward_lines}. Родитель читает \texttt{stdin}
блоками по 4096 байт (константа \fn{kBlockSize}, чтение "--- \fn{read_retry}) и ищет в блоке символы
\texttt{\textbackslash n} функцией \texttt{memchr}. Фрагмент до перевода строки
включительно уходит в текущий канал (\fn{write_all}), после него получатель
меняется: индекс получателя "--- номер строки по модулю \fn{kChildrenCount}. Номер
строки увеличивается только при встрече \texttt{\textbackslash n}, поэтому
строка, разрезанная границей блоков, целиком попадает в один канал:

\begin{lstlisting}[style=console]
block:  "abc\nde"  |  "f\nghi\n"
         pipe1 p2      p2  pipe1
\end{lstlisting}

Так выполняется требование о строках \enquote{произвольной длины}: строка
нигде не хранится целиком, и даже строка, не помещающаяся в оперативную
память, обрабатывается с постоянным расходом памяти. Проверено на строке
длиной 1~ГБ: пиковое потребление памяти родителя около 0{,}9~МБ.

\subsection{Работа дочернего процесса}
Функция \coderef{slave-main}{main} программы \texttt{slave} читает
\texttt{stdin} блоками, удаляет гласные на месте (проверка символа "---
\fn{is_vowel}), сдвигая остальные байты к началу буфера, и записывает результат
в \texttt{stdout} функцией \fn{write_all}, которая повторяет \texttt{write}
при частичной записи. Обработка побайтовая, поэтому место разреза строки
родителем значения не имеет. \texttt{read} возвращает 0, когда все
пишущие концы канала закрыты.

\subsection{Завершение}
Функция \fn{finish_children} закрывает \texttt{pipe1[1]} и \texttt{pipe2[1]},
дети получают EOF и завершаются, после чего для каждого из них вызывается
\fn{wait_child}, которая ожидает процесс через \texttt{waitpid} и разбирает его
статус.
Порядок важен: если ждать детей до закрытия каналов, возникнет взаимная
блокировка "--- ребёнок ждёт конца данных, а родитель ждёт ребёнка.

Все ресурсы (дескрипторы, идентификаторы процессов) хранятся в структуре
\fn{child_t}, где значение $-1$ означает \enquote{не выделено}. Функция
подготовки \fn{start_children} при ошибке просто возвращает код, а
\fn{finish_children} всегда закрывает всё открытое и дожидается всех запущенных
процессов; \coderef{master-main}{main} лишь вызывает эти этапы по очереди. Это позволяет
обойтись без \texttt{goto} и корректно освобождать ресурсы на любом этапе,
например если второй \texttt{fork} не удался, а первый ребёнок уже работает.

\section{Обработка ошибок}
Каждый системный вызов проверяется; вызовы \texttt{read}, \texttt{write},
\texttt{open}, \texttt{waitpid}, прерванные сигналом (\texttt{EINTR}),
повторяются (\fn{read_retry}, \fn{write_all}, \fn{open_output_file},
\fn{wait_child}). Код завершения родителя "--- битовая маска, поэтому по нему видно
все произошедшие ошибки одновременно:

\begin{longtable}{|c|p{0.30\textwidth}|>{\raggedright\arraybackslash}p{0.50\textwidth}|}
\hline
\rowcolor{lightgray}
\textbf{Бит} & \textbf{Константа} & \textbf{Значение} \\
\hline
\endhead
1 & \fn{kErrInput} & ошибка чтения \texttt{stdin}, пустое или слишком
                         длинное имя, один и тот же файл для обоих детей \\
\hline
2 & \fn{kErrSetup} & ошибка \texttt{open}, \texttt{pipe}, \texttt{fork},
                         \texttt{sigaction}, закрыт стандартный дескриптор \\
\hline
4 & \fn{kErrChild1Write} & \texttt{child1} перестал читать \texttt{pipe1} \\
\hline
8 & \fn{kErrChild2Write} & \texttt{child2} перестал читать \texttt{pipe2} \\
\hline
16 & \fn{kErrChild1Status} & \texttt{child1} завершился с ошибкой или по сигналу \\
\hline
32 & \fn{kErrChild2Status} & \texttt{child2} завершился с ошибкой или по сигналу \\*
\hline
\end{longtable}

Максимальное значение 63 не пересекается с кодами 126, 127 и $128+N$, которыми
командная оболочка обозначает ошибки запуска и завершение по сигналу. Дочерний
процесс завершается с кодом 126 (\fn{kChildSetupFailed}), если не удался
\texttt{dup2}, и 127 (\fn{kChildExecFailed}), если не удался \texttt{execv};
сама программа \texttt{slave} возвращает 1 (\fn{kSlaveReadFailed}) при ошибке
чтения и 2 (\fn{kSlaveWriteFailed}) при ошибке записи. Родитель расшифровывает
эти коды в сообщениях (\fn{wait_child}).

Если ребёнок аварийно завершился, запись в его канал возвращает \texttt{EPIPE}:
\fn{forward_lines} сообщает об этом, закрывает канал и продолжает пересылать
строки второму ребёнку.

\section{Демонстрация работы}
Программа разрабатывалась на macOS; трассировка выполнена в контейнере Debian
(ядро Linux 6.12, \texttt{gcc} 14.2, \texttt{strace} 6.13). Входные данные:

\begin{lstlisting}[style=console]
$ cat input.txt
out1.txt
out2.txt
Hello World
Operating Systems
abcdefghij
UNIX pipes
$ ./master < input.txt; echo "exit=$?"
exit=0
$ cat out1.txt
Hll Wrld
bcdfghj
$ cat out2.txt
prtng Systms
NX pps
\end{lstlisting}

\subsection{Трассировка системных вызовов}
\begin{lstlisting}[style=console]
strace -f -o trace.log -e trace=read,write,openat,pipe2,fcntl,clone,execve,\
    dup3,close,wait4,rt_sigaction,exit_group ./master < input.txt
\end{lstlisting}

Фрагмент трассы (строки динамического загрузчика опущены, PID 532 "---
\texttt{master}, 533 и 534 "--- дочерние процессы):

\begin{lstlisting}[style=console]
532   fcntl(0, F_GETFD)                 = 0
532   fcntl(1, F_GETFD)                 = 0
532   fcntl(2, F_GETFD)                 = 0
532   read(0, "o", 1)                   = 1
532   read(0, "u", 1)                   = 1
...
532   read(0, "\n", 1)                  = 1
532   openat(AT_FDCWD, "out1.txt", O_WRONLY|O_CREAT|O_TRUNC|O_CLOEXEC, 0644) = 3
532   openat(AT_FDCWD, "out2.txt", O_WRONLY|O_CREAT|O_TRUNC|O_CLOEXEC, 0644) = 4
532   pipe2([5, 6], 0)                  = 0
532   fcntl(5, F_SETFD, FD_CLOEXEC)     = 0
532   fcntl(6, F_SETFD, FD_CLOEXEC)     = 0
532   pipe2([7, 8], 0)                  = 0
...
532   clone(child_stack=NULL, flags=CLONE_CHILD_CLEARTID|CLONE_CHILD_SETTID|SIGCHLD, ...) = 533
532   close(5)                          = 0
533   dup3(5, 0, 0)                     = 0
532   close(3)                          = 0
533   dup3(3, 1, 0)                     = 1
533   execve("./slave", ["slave"], 0xffffc81fe8c8 /* 7 vars */) = 0
532   clone(...)                        = 534
532   close(7)                          = 0
534   dup3(7, 0, 0)                     = 0
532   close(4)                          = 0
534   dup3(4, 1, 0)                     = 1
532   rt_sigaction(SIGPIPE, {sa_handler=SIG_IGN, sa_mask=[], sa_flags=0}, NULL, 8) = 0
534   execve("./slave", ["slave"], 0xffffc81fe8c8 /* 7 vars */) = 0
532   read(0, "Hello World\nOperating Systems\nab"..., 4096) = 52
532   write(6, "Hello World\n", 12)     = 12
532   write(8, "Operating Systems\n", 18) = 18
532   write(6, "abcdefghij\n", 11)      = 11
532   write(8, "UNIX pipes\n", 11)      = 11
532   read(0, "", 4096)                 = 0
532   close(6)                          = 0
532   close(8)                          = 0
532   wait4(533,  <unfinished ...>
533   read(0, "Hello World\nabcdefghij\n", 4096) = 23
533   write(1, "Hll Wrld\nbcdfghj\n", 17) = 17
534   read(0, "Operating Systems\nUNIX pipes\n", 4096) = 29
534   write(1, "prtng Systms\nNX pps\n", 20) = 20
533   read(0, "", 4096)                 = 0
533   exit_group(0)                     = ?
532   <... wait4 resumed>[{WIFEXITED(s) && WEXITSTATUS(s) == 0}], 0, NULL) = 533
534   read(0, "", 4096)                 = 0
534   exit_group(0)                     = ?
532   wait4(534, [{WIFEXITED(s) && WEXITSTATUS(s) == 0}], 0, NULL) = 534
532   exit_group(0)                     = ?
\end{lstlisting}

По трассе видны все этапы работы:
\begin{itemize}
    \item имена файлов читаются по одному байту (\fn{read_filename}), а строки "---
          одним блоком (\texttt{read(0, ..., 4096) = 52}) в \fn{forward_lines};
    \item glibc реализует \texttt{pipe()} через \texttt{pipe2(fds, 0)},
          \texttt{fork()} "--- через \texttt{clone} с флагом \texttt{SIGCHLD}, а
          \texttt{dup2()} "--- через \texttt{dup3};
    \item в ребёнке (\fn{spawn_slave}) \texttt{dup3(5, 0)} подставляет \texttt{pipe1[0]} на место
          \texttt{stdin}, \texttt{dup3(3, 1)} "--- файл на место \texttt{stdout};
    \item блок из 52 байт разрезается на четыре строки, которые поочерёдно
          уходят в дескрипторы 6 (\texttt{pipe1}) и 8 (\texttt{pipe2});
    \item каждый ребёнок читает обе свои строки одним вызовом \texttt{read}
          и пишет результат одним \texttt{write} в дескриптор 1, то есть в файл;
    \item после \texttt{close(6)} и \texttt{close(8)} дети получают
          \texttt{read(...) = 0} и завершаются, а родитель собирает их статусы
          через \texttt{wait4} (\fn{finish_children}). Строки \texttt{<unfinished ...>} и
          \texttt{<... resumed>} показывают, что процессы выполняются
          параллельно.
\end{itemize}

\subsection{Сводная статистика}
\begin{lstlisting}[style=console]
strace -f -c ./master < input.txt
\end{lstlisting}

\begin{lstlisting}[style=console]
% time     seconds  usecs/call     calls    errors syscall
------ ----------- ----------- --------- --------- ----------------
 70.17    0.001508         754         2           wait4
  9.35    0.000201           9        21           mmap
  7.21    0.000155          12        12           mprotect
  3.44    0.000074           9         8           munmap
  2.14    0.000046           1        27           read
  2.14    0.000046           7         6           write
  ...
  0.00    0.000000           0         4           dup3
  0.00    0.000000           0        11           fcntl
  0.00    0.000000           0         2           pipe2
  0.00    0.000000           0         1           rt_sigaction
  0.00    0.000000           0         2           clone
  0.00    0.000000           0         3           execve
------ ----------- ----------- --------- --------- ----------------
100.00    0.002149          13       156         6 total
\end{lstlisting}

Основное время уходит на \texttt{wait4}: родитель блокируется в нём, пока дети
обрабатывают данные. Три вызова \texttt{execve} "--- запуск \texttt{master} и
двух экземпляров \texttt{slave}; \texttt{mmap}, \texttt{mprotect} и
\texttt{munmap} выполняет динамический загрузчик при старте каждого из трёх
процессов.

\section{Итог работы}
В ходе лабораторной работы была реализована программа из двух
процессов-обработчиков и управляющего процесса, взаимодействующих через
неименованные каналы. Получены практические навыки работы с системными вызовами
\texttt{fork}, \texttt{execv}, \texttt{dup2}, \texttt{pipe}, \texttt{waitpid}
и \texttt{sigaction}.

Основные выводы:
\begin{itemize}
    \item каждый лишний открытый конец канала "--- потенциальная взаимоблокировка,
          поэтому их удобно закрывать автоматически флагом \texttt{FD\_CLOEXEC};
    \item потоковая обработка блоками фиксированного размера позволяет работать
          со строками любой длины при постоянном расходе памяти;
    \item игнорирование \texttt{SIGPIPE} превращает аварию дочернего процесса в
          обрабатываемую ошибку \texttt{EPIPE};
    \item порядок завершения важен: сначала закрыть пишущие концы каналов, затем
          ожидать дочерние процессы.
\end{itemize}

\newpage
\phantomsection
\addcontentsline{toc}{section}{Приложение А. Исходный код программы}
\begin{center}
    {\large\bfseries ПРИЛОЖЕНИЕ А}\\[0.2cm]
    {\bfseries Исходный код программы}
\end{center}

Исходный код также доступен в
\MYhref{https://github.com/evdakCherpak/os_mai_labs}{репозитории на GitHub}.

\phantomsection
\addcontentsline{toc}{subsection}{А.1. master.c}
\subsection*{А.1. master.c}
\lstinputlisting[name=master.c]{../src/master.c}

\phantomsection
\addcontentsline{toc}{subsection}{А.2. slave.c}
\subsection*{А.2. slave.c}
\lstinputlisting[name=slave.c]{../src/slave.c}

\phantomsection
\addcontentsline{toc}{subsection}{А.3. common.h}
\subsection*{А.3. common.h}
\lstinputlisting[name=common.h]{../src/common.h}

\phantomsection
\addcontentsline{toc}{subsection}{А.4. common.c}
\subsection*{А.4. common.c}
\lstinputlisting[name=common.c]{../src/common.c}

\phantomsection
\addcontentsline{toc}{subsection}{А.5. Makefile}
\subsection*{А.5. Makefile}
% escapechar=@ from log-style would break on $@ in the Makefile
\lstinputlisting[name=Makefile,language=make,escapechar={}]{../Makefile}

\end{document}
